\paragraph{2.1. Сортирока вставкой}
\subparagraph{2.1.1}
Картинка* не знаю как вставить** да и готовить её надо, лень

\subparagraph{2.1.2}
Перепишите процедуру Insertion-Sort для сортировки в невозрастающем порядке вместо неубывающего. \\
\\
\begin{tabular}{ll}
Insertion-Sort(A) & \\
\quad for j = 2 to A.length & \\
\quad \quad key = A[j] & \\
\quad \quad i = j - 1 & \\
\quad \quad while i > 0 and A[i] < key & \\
\quad \quad \quad A[i+1] = A[i] & \\
\quad \quad \quad i = i - 1 & \\
\quad \quad A[i+1] = key &
\end{tabular}

\subparagraph{2.1.3}
Рассмотрите {\bfseries задачу поиска.}

{\bfseries Вход.} Последовательность из $n$ чисел $A = \left\langle a_1, a_2, \ldots, a_n \right\rangle$ и значение $\upsilon$.

{\bfseries Выход.} Индекс $i$, такой, что $\upsilon = A[i]$, или специальное значение NIL, если $\upsilon$ в $A$ отсутствует.

Составьте псевдокод {\bfseries линейного поиска}, при работе которого выполняется сканирование последовательности в поисках значения $\upsilon$. Докажите корректность алгоритма с помощью инварианта цикла. Убедитесь, что выбранный инвариант цикла удовлетворяет трём необходимым условиям. \\

\begin{tabular}{ll}
Linear-Search(A, $\upsilon$) & \\
\quad for $i = 1$ to A.length & \\
\quad \quad if $A[i]$ == $\upsilon$ & \\
\quad \quad \quad return $i$ & \\
\quad return NILL & 
\end{tabular} \\
\\
\textbf{Инвариант цикла:} \\
В начале каждой итерации цикла for подмассив $A[1 \ldots i-1]$ не содержит значение $\upsilon$.

\textbf{Инициализация:} \\
При $i = 1$ подмассив $A[1 \ldots 0]$ пуст, поэтому он не содержит $\upsilon$.

\textbf{Сохранение:} \\
Предположим, что инвариант верен перед началом итерации с текущим $i$. Проверяем элемент $A[i]$:
\begin{itemize}
    \item Если $A[i] = \upsilon$, алгоритм завершается, возвращая $i$. Перед завершением инвариант выполнен, так как $A[1 \ldots i-1]$ не содержит $\upsilon$.
    \item Если $A[i] \neq \upsilon$, то после увеличения $i$ на 1 подмассив $A[1 \ldots i-1]$ (в следующей итерации) по-прежнему не содержит $\upsilon$, так как к предыдущим элементам добавился $A[i]$, который не равен $\upsilon$.
\end{itemize}

\textbf{Завершение:} \\
Цикл завершается в одном из двух случаев:
\begin{enumerate}
    \item Найден индекс $i$, такой что $A[i] = \upsilon$. Тогда инвариант гарантирует, что в $A[1 \ldots i-1]$ значения $\upsilon$ не было.
    \item Перебраны все элементы ($i = n+1$). Тогда инвариант означает, что ни один элемент массива не равен $\upsilon$, и алгоритм возвращает $NIL$.
\end{enumerate}

\subparagraph{2.1.4}
Рассмотрим задачу сложения двух n-битовых двоичных целых чисел, хранящихся в n-элементных массивах A и B. Сумму этих двух чисел необходимо занести в двоичной форме в (n+1)-элементном массиве C. Приведите строгую формулировку задачи и составьте псевдокод для сложения этих двух чисел. \\

\textbf{Формулировка:}

Даны два массива $A[1..n]$ и $B[1..n]$, каждый из которых содержит двоичные цифры (0 или 1), представляющие два $n$-битовых двоичных числа:
\[
A = a_1 a_2 \ldots a_n, \quad B = b_1 b_2 \ldots b_n,
\]
где $a_1$ и $b_1$ — младшие биты. Требуется вычислить их сумму $C = A + B$, которая является $(n+1)$-битовым двоичным числом, и записать результат в массив $C[1..n+1]$. \\

\textbf{Псевдокод:} \\
\begin{tabular}{ll}
Bitwise-Addition($A, B, C$) & \\
\quad $carry = 0$ & \\
\quad for $i = 1$ to A.length & \\
\quad \quad $sum = A[i] + B[i] + carry$ & \\
\quad \quad $C[i] = sum\bmod{2}$ & \\
\quad \quad $carry = \lfloor sum/2 \rfloor$ & \\
\quad $C[i+1] = carry$ & \\
\end{tabular} \\